% Part: first-order-logic
% Chapter: models-theories
% Section: theories

\documentclass[../../../include/open-logic-section]{subfiles}

\begin{document}

\olfileid{fol}{mat}{the}

\olsection{Voorbeelden van eerste-ordetheorieën}

\begin{ex}
We beginnen met de theorie van strikte lineaire ordeningen. Daarvoor
gebruiken we de taal~$\Lang L_<$. De volgende verzameling zinnen vormt
het axiomasysteem:
\begin{align*}
\{\quad & \lforall[x][\lnot x < x], \\
& \lforall[x][\lforall[y][((x < y \lor y <
    x) \lor x = y)]], \\
& \lforall[x][\lforall[y][\lforall[z][((x < y
      \land y < z) \lif x < z)]]] \quad \}
\end{align*}
Samen leggen deze axioma's precies vast welke !!{structure}s we
bedoelen. Iedere strikte lineaire ordening is een model van dit
axiomasysteem. En omgekeerd: als $R$ een lineaire ordening op een
verzameling~$X$ is, dan is de !!{structure}~$\Struct M$ met
$\Domain M = X$ en $\Assign{<}{M} = R$ een model van deze theorie.
\end{ex}

\begin{ex}
Voor de theorie van groepen gebruiken we een taal met $\Obj 1$ (een
!!{constant}) en $\cdot$ (een tweeplaatsige !!{function}). De theorie
heeft de volgende axioma's:
\begin{align*}
& \lforall[x][\eq[(x \cdot \Obj 1)][x]]\\
& \lforall[x][\lforall[y][\lforall[z][\eq[(x \cdot (y \cdot z))][((x
          \cdot y) \cdot z)]]]]\\
& \lforall[x][\lexists[y][\eq[(x \cdot y)][\Obj 1]]]
\end{align*}
\end{ex}

\begin{ex}
De Peano-rekenkunde---de theorie over de natuurlijke getallen---krijgt
de volgende axioma's in de taal van de rekenkunde~$\Lang L_A$.
\begin{align*}
& \lforall[x][\lforall[y][(\eq[x'][y'] \lif \eq[x][y])]]\\
& \lforall[x][\eq/[\Obj 0][x']]\\
& \lforall[x][\eq[(x + \Obj 0)][x]]\\
& \lforall[x][\lforall[y][\eq[(x + y')][(x + y)']]]\\
& \lforall[x][\eq[(x \times \Obj 0)][\Obj 0]]\\
& \lforall[x][\lforall[y][\eq[(x \times y')][((x \times y) + x)]]]\\
& \lforall[x][\lforall[y][(x < y \liff \lexists[z][\eq[(z' + x)][y])]]]\\
\intertext{en alle zinnen van de volgende vorm}
& (!A(\Obj 0) \land \lforall[x][(!A(x) \lif !A(x'))]) \lif \lforall[x][!A(x)]
\end{align*}
Van die laatste vorm zijn er oneindig veel zinnen. Het axiomasysteem
is dus ook oneindig. Het vaste patroon dat voor iedere formule zo'n
axioma oplevert, heet het \emph{inductieschema}. (Het volledige
inductieschema is eigenlijk iets ingewikkelder dan we hier laten
zien.)

Het laatste axioma is dus een \emph{expliciete definitie} van~$<$: het
legt vast wat dit teken betekent.
\end{ex}

\begin{ex}
De theorie van zuivere verzamelingen is belangrijk voor de grondslagen
en de filosofie van de wiskunde. Een verzameling is zuiver als al haar
!!{element}s zelf ook zuivere verzamelingen zijn. De lege verzameling
is daarom zuiver. Maar zodra een verzameling iets als !!a{element}
heeft dat zelf geen verzameling is, is ze niet zuiver. Je kunt het ook
zo zeggen: je begint met de lege verzameling en bouwt daar alleen
andere verzamelingen uit. Er komen geen ``oerelementen'' in voor, dus
geen objecten die zelf geen verzameling zijn.

Je kunt de volgende zinnen als axioma's voor een theorie van zuivere
verzamelingen nemen:
\begin{align*}
& \lexists[x][\lnot \lexists[y][y \in x]]\\
& \lforall[x][\lforall[y][(\lforall[z](z \in x \liff z \in y) \lif 
\eq[x][y])]]\\
& \lforall[x][\lforall[y][\lexists[z][\lforall[u][(u \in z \liff
(\eq[u][x] \lor \eq[u][y]))]]]]\\
& \lforall[x][\lexists[y][\lforall[z][(z \in y \liff \lexists[u][(z \in
        u \land u \in x)])]]]\\
\intertext{en alle zinnen van de volgende vorm} &
\lexists[x][\lforall[y][(y \in x \liff !A(y))]]
\end{align*}
Het eerste axioma zegt dat er een verzameling zonder !!{element}s is:
$\emptyset$ bestaat. Het tweede zegt dat twee verzamelingen met precies
dezelfde elementen gelijk zijn; dit heet extensionaliteit. Het derde
zegt dat voor iedere twee verzamelingen $X$ en $Y$ ook de verzameling
$\{X, Y\}$ bestaat. Het vierde zegt dat voor iedere verzameling $X$ de
verzameling $\cup X$ bestaat. Die $\cup X$ krijg je door alle elementen
van de verzamelingen in $X$ bij elkaar te nemen.

De !!{sentence}s van de laatste soort heten samen het
\emph{naïeve comprehensieschema}. Voor iedere formule $!A(x)$ zegt dit
patroon dat de verzameling $\Setabs{x}{!A(x)}$ bestaat: de verzameling
van precies de dingen waarvoor die formule geldt. Op het eerste gezicht
lijkt dat waar, nuttig en misschien zelfs noodzakelijk. Maar juist
daarom heet het ``naïef'': met dit schema kan de theorie geen model
hebben. Neem voor $!A(y)$ namelijk $\lnot y \in y$. Dan krijg je de
!!{sentence}
\[
\lexists[x][\lforall[y][(y \in x \liff \lnot y \in y)]]
\]
en die !!{sentence} wordt in geen enkele !!{structure} vervuld.
\end{ex}

\begin{ex}
\emph{Mereologie} is de leer van delen en gehelen. De relatie
\emph{is een deel van} staat daarin centraal. Net zoals er
verschillende theorieën over verzamelingen zijn, zijn er ook
verschillende theorieën over deze deel-geheelrelatie. Hun axioma's
leggen verschillende, soms botsende, opvattingen ervan vast.

De taal van de mereologie heeft maar één tweeplaatsig
predicaatsymbool:~$\Obj P$. De formule $\Atom{\Obj P}{x, y}$
``betekent'' dat $x$ een deel van~$y$ is. Als we dit symbool zo lezen,
noemen we !!a{structure} voor deze taal een
\emph{deel-geheelstructuur}. Niet iedere structuur met één
tweeplaatsig predicaat verdient die naam. De relatie
$\Assign{\Obj P}{M}$ moet eerst aan bepaalde voorwaarden voldoen. Die
voorwaarden kunnen we als axioma's opschrijven. Zo moet ``is een deel
van'' een partiële ordening zijn. Ieder object is een deel van
zichzelf; dat is dan een \emph{oneigenlijk} deel. Twee verschillende
objecten kunnen niet elk een deel van het andere zijn. En als iets een
deel is van een deel van een object, dan is het ook een deel van dat
object zelf. In deze betekenis lijkt ``is een deel van'' dus op ``is
een deelverzameling van''. Het lijkt niet op ``is een element van'',
want die laatste relatie is niet reflexief en ook niet transitief.
\begin{align*}
& \lforall[x][\Atom{\Obj P}{x,x}] \\
& \lforall[x][\lforall[y][((\Part{x}{y} \land \Part{y}{x})
      \lif \eq[x][y])]] \\
& \lforall[x][\lforall[y][\lforall[z][((\Part{x}{y} \land
        \Part{y}{z}) \lif \Part{x}{z})]]]\\
\intertext{Bovendien hebben elke twee objecten een mereologische som: een object
  dat beide als deel heeft en in dat opzicht zo klein mogelijk is.}  &
\lforall[x][\lforall[y][\lexists[z][\lforall[u][(\Part{z}{u} \liff
        (\Part{x}{u} \land \Part{y}{u}))]]]]
\end{align*}
Dit zijn maar enkele basisregels voor delen en gehelen die metafysici
bestuderen. Andere regels zijn al snel lastig precies op te schrijven
als je niet eerst een paar extra relaties definieert. De meeste
metafysici die mereologie bestuderen, vinden bijvoorbeeld ook het
volgende geldig. Stel dat een object~$x$ een strikt deel~$y$ heeft,
dus een deel dat niet het hele object is. Dan heeft het ook een deel~$z$
dat geen enkel deel met~$y$ gemeen heeft. De fusie van $y$ en $z$---hun
mereologische som---is dan precies~$x$.
\end{ex}

\end{document}
